\chapter{Aide-mémoire}

\begin{synthese}[Mode d'emploi]
Cette annexe rassemble, sous forme de fiches de référence denses, l'essentiel de
la syntaxe vue dans le fascicule : \textbf{PHP}, \textbf{SQL}, l'extension
\textbf{mysqli} (PHP + MySQL) et le langage \textbf{C}. Une dernière fiche met en
correspondance PHP et C. À consulter pendant les TP et les révisions.
\end{synthese}

\section{Aide-mémoire PHP}

\subsection{Balises, commentaires et instructions}

\begin{syntaxe}
\code{<?php}\quad \textit{instructions\dots}\quad \code{?>} \qquad — fichier en
\code{.php}, chaque instruction se termine par \code{;}.
\end{syntaxe}

\begin{codephp}
<?php
  // commentaire sur une ligne
  # autre commentaire sur une ligne
  /* commentaire
     sur plusieurs lignes */
  echo "Bonjour OnBuch";   // une instruction = un point-virgule
?>
\end{codephp}

\subsection{Variables et constantes}

Une variable commence \textbf{toujours} par \code{\$}. PHP est typé
dynamiquement : le type dépend de la valeur affectée.

\begin{codephp}
<?php
  $nom    = "Awa";              // string  (chaîne)
  $age    = 17;                 // int     (entier)
  $moy    = 14.5;               // float   (décimal)
  $admis  = true;               // bool    (true / false)
  define("VILLE", "Maroua");    // constante (sans $)
  const TVA = 0.1925;           // autre forme de constante
?>
\end{codephp}

\begin{center}\small
\begin{tabular}{@{}ll@{}}
\toprule
\textbf{Type} & \textbf{Exemple de valeur}\\
\midrule
\code{int}    & \code{42}, \code{-7}, \code{0}\\
\code{float}  & \code{14.5}, \code{3.14}\\
\code{string} & \code{"Awa"}, \code{'Lycée'}\\
\code{bool}   & \code{true}, \code{false}\\
\code{array}  & \code{[1, 2, 3]}\\
\code{NULL}   & \code{null}\\
\bottomrule
\end{tabular}
\end{center}

\subsection{Afficher : \texttt{echo} et \texttt{print}}

\begin{itemize}
  \item \code{echo} affiche une ou plusieurs valeurs séparées par des virgules ;
  \item \code{print} affiche une seule valeur et renvoie \code{1} ;
  \item \code{.} concatène (colle) deux chaînes ; \code{.=} ajoute à la fin.
\end{itemize}

\begin{codephp}
<?php
  $nom = "Awa"; $age = 17;
  echo "Élève : $nom",  " (", $age, " ans)";   // virgules
  echo "Élève : " . $nom . " (" . $age . ")";  // concaténation
  echo "Bonjour {$nom} !";                      // accolades dans "..."
  echo 'Texte brut : $nom';                     // '...' n'interprète pas $
?>
\end{codephp}

\subsection{Opérateurs}

\begin{center}\small
\begin{tabular}{@{}lll@{}}
\toprule
\textbf{Catégorie} & \textbf{Opérateurs} & \textbf{Exemple}\\
\midrule
Arithmétiques & \code{+ - * / \% **} & \code{\$x \% 2} (reste)\\
Affectation   & \code{= += -= *= /= .=} & \code{\$s .= "!"}\\
Comparaison   & \code{== != < > <= >=} & \code{\$a <= \$b}\\
Identité      & \code{=== !==} & \code{\$a === \$b} (valeur + type)\\
Logiques      & \code{\&\& || !} & \code{\$x>0 \&\& \$x<10}\\
Incrémentation & \code{++ -{}-} & \code{\$i++}\\
Ternaire      & \code{? :} & \code{\$m>=10 ? "Admis":"Échec"}\\
\bottomrule
\end{tabular}
\end{center}

\subsection{Structures de contrôle}

\begin{codephp}
<?php
  // if / elseif / else
  if ($moy >= 16)      { echo "Très bien"; }
  elseif ($moy >= 10)  { echo "Admis"; }
  else                 { echo "Échec"; }

  // switch
  switch ($jour) {
    case 1: echo "Lundi"; break;
    case 2: echo "Mardi"; break;
    default: echo "Autre";
  }
?>
\end{codephp}

\begin{codephp}
<?php
  // for : compteur connu
  for ($i = 1; $i <= 5; $i++) { echo $i; }

  // while : tant que la condition est vraie
  $n = 1;
  while ($n <= 5) { echo $n; $n++; }

  // do...while : exécute au moins une fois
  $n = 1;
  do { echo $n; $n++; } while ($n <= 5);

  // foreach : parcourir un tableau
  foreach ($notes as $note) { echo $note; }
  foreach ($eleve as $cle => $valeur) { echo "$cle = $valeur"; }
?>
\end{codephp}

\begin{remarque}
\code{break} sort de la boucle (ou du \code{switch}) ; \code{continue} passe
directement à l'itération suivante.
\end{remarque}

\subsection{Tableaux}

\begin{codephp}
<?php
  // tableau indexé (clés 0, 1, 2...)
  $notes = [12, 15, 9, 18];
  echo $notes[0];           // 12
  $notes[] = 11;            // ajoute en fin

  // tableau associatif (clé => valeur)
  $eleve = ["nom" => "Awa", "age" => 17, "ville" => "Maroua"];
  echo $eleve["nom"];       // Awa
?>
\end{codephp}

\begin{center}\small
\begin{tabular}{@{}ll@{}}
\toprule
\textbf{Fonction} & \textbf{Rôle}\\
\midrule
\code{count(\$t)}            & nombre d'éléments\\
\code{sort(\$t)}             & trie par valeurs croissantes\\
\code{rsort(\$t)}            & trie par valeurs décroissantes\\
\code{in\_array(\$v, \$t)}   & \code{true} si \code{\$v} est dans \code{\$t}\\
\code{array\_push(\$t, \$v)} & ajoute \code{\$v} en fin\\
\code{implode(",", \$t)}     & tableau $\rightarrow$ chaîne « 12,15,9 »\\
\code{explode(",", \$s)}     & chaîne $\rightarrow$ tableau\\
\bottomrule
\end{tabular}
\end{center}

\subsection{Fonctions}

\begin{syntaxe}
\code{function nom(\$param1, \$param2 = défaut) \{ \dots return \$valeur; \}}
\end{syntaxe}

\begin{codephp}
<?php
  function moyenne($a, $b, $c) {
    return ($a + $b + $c) / 3;   // return renvoie le résultat
  }
  function bonjour($nom = "élève") {  // valeur par défaut
    return "Bonjour $nom";
  }
  echo moyenne(12, 15, 9);   // 12
  echo bonjour();            // Bonjour élève
?>
\end{codephp}

\subsection{Superglobales (formulaires \& serveur)}

\begin{center}\small
\begin{tabular}{@{}ll@{}}
\toprule
\textbf{Superglobale} & \textbf{Contenu}\\
\midrule
\code{\$\_GET}     & données passées dans l'URL (\code{?id=3})\\
\code{\$\_POST}    & données d'un formulaire (\code{method="post"})\\
\code{\$\_REQUEST} & \code{\$\_GET} + \code{\$\_POST} réunis\\
\code{\$\_FILES}   & fichiers téléversés (upload)\\
\code{\$\_SERVER}  & infos serveur (\code{REQUEST\_METHOD}, \dots)\\
\code{\$\_SESSION} & variables de session\\
\bottomrule
\end{tabular}
\end{center}

\begin{codephp}
<?php
  // récupérer un champ de formulaire envoyé en POST
  $nom = $_POST["nom"];
  $age = $_GET["age"];

  if ($_SERVER["REQUEST_METHOD"] == "POST") {
    echo "Formulaire envoyé : " . $_POST["nom"];
  }
?>
\end{codephp}

\begin{attention}
Ne jamais faire confiance aux données reçues. Vérifie l'existence avec
\code{isset(\$\_POST["nom"])} avant de les utiliser.
\end{attention}

\subsection{Fonctions utiles}

\begin{center}\small
\begin{tabular}{@{}lll@{}}
\toprule
\textbf{Fonction} & \textbf{Rôle} & \textbf{Exemple $\rightarrow$ résultat}\\
\midrule
\code{strlen(\$s)}        & longueur d'une chaîne & \code{strlen("Awa")} $\rightarrow$ 3\\
\code{strtoupper(\$s)}    & majuscules            & \code{"awa"} $\rightarrow$ \code{AWA}\\
\code{strtolower(\$s)}    & minuscules            & \code{"AWA"} $\rightarrow$ \code{awa}\\
\code{trim(\$s)}          & enlève les espaces    & \code{" x "} $\rightarrow$ \code{x}\\
\code{round(\$n, 2)}      & arrondit              & \code{round(14.567,2)} $\rightarrow$ 14.57\\
\code{abs(\$n)}           & valeur absolue        & \code{abs(-5)} $\rightarrow$ 5\\
\code{rand(1, 6)}         & entier aléatoire      & \code{rand(1,6)} $\rightarrow$ 4\\
\code{date("d/m/Y")}      & date du jour          & $\rightarrow$ \code{26/06/2026}\\
\code{isset(\$v)}         & \code{\$v} existe ?   & \code{true / false}\\
\bottomrule
\end{tabular}
\end{center}

\section{Aide-mémoire SQL}

\subsection{Types de données courants}

\begin{center}\small
\begin{tabular}{@{}ll@{}}
\toprule
\textbf{Type} & \textbf{Usage}\\
\midrule
\code{INT}          & nombre entier (âge, identifiant)\\
\code{FLOAT}        & nombre décimal (moyenne, prix)\\
\code{VARCHAR(n)}   & chaîne de longueur variable, max \code{n}\\
\code{TEXT}         & texte long\\
\code{DATE}         & date au format \code{AAAA-MM-JJ}\\
\code{BOOLEAN}      & vrai / faux (0 ou 1)\\
\bottomrule
\end{tabular}
\end{center}

\subsection{Créer une table}

\begin{codesql}
CREATE TABLE eleves (
  id        INT PRIMARY KEY AUTO_INCREMENT,
  nom       VARCHAR(50) NOT NULL,
  age       INT,
  moyenne   FLOAT,
  ville     VARCHAR(40),
  inscrit   DATE
);
\end{codesql}

\begin{remarque}
La \textbf{clé primaire} (\code{PRIMARY KEY}) identifie chaque ligne de façon unique.
\code{AUTO\_INCREMENT} fait que la base attribue elle-même les valeurs 1, 2, 3, \dots
\end{remarque}

\subsection{Insérer des données}

\begin{codesql}
INSERT INTO eleves (nom, age, moyenne, ville, inscrit)
VALUES ('Awa', 17, 14.5, 'Maroua', '2026-09-01');

INSERT INTO eleves (nom, age) VALUES ('Bello', 18);
\end{codesql}

\subsection{Lire : SELECT}

\begin{codesql}
SELECT * FROM eleves;                       -- toutes les colonnes
SELECT nom, moyenne FROM eleves;            -- colonnes choisies

SELECT * FROM eleves
WHERE moyenne >= 10                         -- filtre (condition)
ORDER BY moyenne DESC;                      -- tri décroissant

SELECT * FROM eleves WHERE ville = 'Maroua' AND age > 16;
SELECT COUNT(*) FROM eleves;                -- nombre de lignes
\end{codesql}

\begin{center}\small
\begin{tabular}{@{}ll@{}}
\toprule
\textbf{Clause} & \textbf{Rôle}\\
\midrule
\code{WHERE}    & condition de sélection\\
\code{ORDER BY} & tri (\code{ASC} croissant, \code{DESC} décroissant)\\
\code{AND OR}   & combiner des conditions\\
\code{LIKE}     & motif texte (\code{'A\%'} = commence par A)\\
\code{LIMIT n}  & limiter à \code{n} lignes\\
\bottomrule
\end{tabular}
\end{center}

\subsection{Modifier et supprimer}

\begin{codesql}
UPDATE eleves SET moyenne = 16 WHERE nom = 'Awa';

DELETE FROM eleves WHERE id = 5;
\end{codesql}

\begin{attention}
Un \code{UPDATE} ou un \code{DELETE} \textbf{sans} \code{WHERE} modifie ou supprime
\textbf{toutes} les lignes de la table. Toujours préciser la condition !
\end{attention}

\section{Aide-mémoire mysqli (PHP + MySQL)}

\subsection{Fonctions essentielles}

\begin{center}\small
\begin{tabular}{@{}ll@{}}
\toprule
\textbf{Fonction} & \textbf{Rôle}\\
\midrule
\code{mysqli\_connect(h,u,p,bd)} & se connecter au serveur MySQL\\
\code{mysqli\_select\_db(\$c,bd)} & choisir la base de données\\
\code{mysqli\_query(\$c, \$sql)}  & exécuter une requête SQL\\
\code{mysqli\_fetch\_assoc(\$r)}  & 1 ligne en tableau associatif\\
\code{mysqli\_fetch\_row(\$r)}    & 1 ligne en tableau indexé\\
\code{mysqli\_fetch\_array(\$r)}  & 1 ligne (indexé + associatif)\\
\code{mysqli\_num\_rows(\$r)}     & nombre de lignes du résultat\\
\code{mysqli\_close(\$c)}         & fermer la connexion\\
\code{mysqli\_connect\_error()}   & message d'erreur de connexion\\
\bottomrule
\end{tabular}
\end{center}

\subsection{Squelette de connexion}

\begin{codephp}
<?php
  $serveur = "localhost";
  $user    = "root";
  $passe   = "";
  $base    = "onbuch";

  // connexion (die() arrête le script en cas d'échec)
  $conn = mysqli_connect($serveur, $user, $passe, $base)
          or die("Connexion impossible : " . mysqli_connect_error());
?>
\end{codephp}

\subsection{Lire et afficher (boucle)}

\begin{codephp}
<?php
  $sql = "SELECT nom, moyenne FROM eleves ORDER BY moyenne DESC";
  $res = mysqli_query($conn, $sql);

  if (mysqli_num_rows($res) > 0) {
    // parcourir chaque ligne du résultat
    while ($ligne = mysqli_fetch_assoc($res)) {
      echo $ligne["nom"] . " : " . $ligne["moyenne"] . "<br>";
    }
  } else {
    echo "Aucun élève trouvé.";
  }

  mysqli_close($conn);   // toujours fermer à la fin
?>
\end{codephp}

\begin{remarque}
\code{mysqli\_fetch\_assoc} renvoie un tableau dont les clés sont les \emph{noms
de colonnes} (\code{\$ligne["nom"]}). \code{mysqli\_fetch\_row} renvoie des clés
numériques (\code{\$ligne[0]}).
\end{remarque}

\subsection{Insérer depuis un formulaire}

\begin{codephp}
<?php
  $nom = $_POST["nom"];
  $age = $_POST["age"];
  $sql = "INSERT INTO eleves (nom, age) VALUES ('$nom', $age)";
  if (mysqli_query($conn, $sql)) {
    echo "Élève ajouté !";
  } else {
    echo "Erreur : " . mysqli_error($conn);
  }
?>
\end{codephp}

\section{Aide-mémoire langage C}

\subsection{Structure d'un programme}

\begin{codec}
#include <stdio.h>      /* bibliothèque entrées / sorties */
#include <stdlib.h>

int main() {
    printf("Bonjour OnBuch\n");
    return 0;           /* 0 = le programme s'est bien terminé */
}
\end{codec}

\subsection{Types et déclarations}

\begin{center}\small
\begin{tabular}{@{}lll@{}}
\toprule
\textbf{Type} & \textbf{Contenu} & \textbf{Format printf/scanf}\\
\midrule
\code{int}    & entier        & \code{\%d}\\
\code{float}  & décimal       & \code{\%f}\\
\code{double} & décimal précis & \code{\%lf}\\
\code{char}   & 1 caractère   & \code{\%c}\\
\code{char[]} & chaîne        & \code{\%s}\\
\bottomrule
\end{tabular}
\end{center}

\begin{codec}
int age = 17;             /* déclaration + initialisation */
float moyenne = 14.5;
char lettre = 'A';
char nom[30] = "Awa";     /* chaîne = tableau de char */
\end{codec}

\subsection{Afficher et lire : printf / scanf}

\begin{codec}
int age;
float moy;
printf("Entre ton age : ");
scanf("%d", &age);              /* & = adresse de la variable */
printf("Entre ta moyenne : ");
scanf("%f", &moy);
printf("Age : %d, moyenne : %.2f\n", age, moy);
\end{codec}

\begin{attention}
Dans \code{scanf}, on met \code{\&} devant le nom de la variable (sauf pour une
chaîne \code{char[]}). \code{\textbackslash n} = retour à la ligne,
\code{\%.2f} = 2 décimales.
\end{attention}

\subsection{Opérateurs}

\begin{center}\small
\begin{tabular}{@{}ll@{}}
\toprule
\textbf{Catégorie} & \textbf{Opérateurs}\\
\midrule
Arithmétiques & \code{+ - * / \%}\\
Comparaison   & \code{== != < > <= >=}\\
Logiques      & \code{\&\& || !}\\
Affectation   & \code{= += -= *= /=}\\
Incrémentation & \code{++ -{}-}\\
\bottomrule
\end{tabular}
\end{center}

\subsection{Structures de contrôle}

\begin{codec}
/* if / else if / else */
if (moy >= 16)       printf("Tres bien");
else if (moy >= 10)  printf("Admis");
else                 printf("Echec");

/* switch */
switch (jour) {
    case 1: printf("Lundi"); break;
    case 2: printf("Mardi"); break;
    default: printf("Autre");
}
\end{codec}

\begin{codec}
/* for */
for (int i = 1; i <= 5; i++) printf("%d ", i);

/* while */
int n = 1;
while (n <= 5) { printf("%d ", n); n++; }

/* do...while : exécute au moins une fois */
n = 1;
do { printf("%d ", n); n++; } while (n <= 5);
\end{codec}

\subsection{Fonctions}

\begin{syntaxe}
\code{type nom(type param1, type param2) \{ \dots return valeur; \}}
\end{syntaxe}

\begin{codec}
/* fonction qui renvoie un float */
float moyenne(int a, int b, int c) {
    return (a + b + c) / 3.0;
}

/* fonction sans retour : void */
void saluer(char nom[]) {
    printf("Bonjour %s\n", nom);
}

int main() {
    printf("%.2f\n", moyenne(12, 15, 9));
    saluer("Awa");
    return 0;
}
\end{codec}

\subsection{Tableaux et structures}

\begin{codec}
/* tableau de 4 entiers */
int notes[4] = {12, 15, 9, 18};
printf("%d", notes[0]);          /* 12 */

for (int i = 0; i < 4; i++)
    printf("%d ", notes[i]);

/* structure (struct) : regrouper plusieurs champs */
struct Eleve {
    char nom[30];
    int  age;
    float moyenne;
};

struct Eleve e1 = {"Awa", 17, 14.5};
printf("%s a %d ans\n", e1.nom, e1.age);   /* accès par . */
\end{codec}

\section{Correspondance PHP \texorpdfstring{$\leftrightarrow$}{<->} C}

Les deux langages partagent une syntaxe proche (héritée du C). Voici les
différences à connaître pour passer de l'un à l'autre.

\begin{center}\small
\begin{tabular}{@{}lll@{}}
\toprule
\textbf{Élément} & \textbf{PHP} & \textbf{C}\\
\midrule
Variable     & \code{\$age = 17;}            & \code{int age = 17;}\\
Décimal      & \code{\$m = 14.5;}            & \code{float m = 14.5;}\\
Chaîne       & \code{\$nom = "Awa";}         & \code{char nom[] = "Awa";}\\
Affichage    & \code{echo \$age;}            & \code{printf("\%d", age);}\\
Saisie       & \code{\$\_POST["x"]}          & \code{scanf("\%d", \&x);}\\
Condition    & \code{if (\$x > 0) \{ \}}     & \code{if (x > 0) \{ \}}\\
Boucle for   & \code{for(\$i=0;\$i<5;\$i++)} & \code{for(int i=0;i<5;i++)}\\
Fonction     & \code{function f()\{return;\}} & \code{int f()\{return 0;\}}\\
Commentaire  & \code{// \dots} ou \code{\# \dots} & \code{// \dots} ou \code{/* \dots */}\\
Tableau      & \code{[12, 15, 9]}            & \code{\{12, 15, 9\}}\\
Fin d'instr. & \code{;}                      & \code{;}\\
\bottomrule
\end{tabular}
\end{center}

\begin{remarque}
Différence majeure : en \textbf{PHP}, on ne déclare pas le type des variables (typage
dynamique) et tout commence par \code{\$} ; en \textbf{C}, chaque variable a un type
fixe déclaré à l'avance (typage statique) et il faut compiler le programme avant de
l'exécuter.
\end{remarque}

\begin{aretenir}
Garde ces fiches à portée de main : \textbf{PHP} pour le web dynamique côté serveur,
\textbf{SQL} pour interroger la base, \textbf{mysqli} pour relier les deux, et
\textbf{C} pour la programmation procédurale. La logique (conditions, boucles,
fonctions) est la même partout : seule la syntaxe change.
\end{aretenir}
